\documentclass[12pt]{article}
\usepackage[utf8]{inputenc}
\usepackage[T1]{fontenc}
\usepackage{hyperref}
\title{Tal Language Resources for K-12 Curriculum Development}
\author{Compiled by Claude Sonnet 5 (claude-sonnet-5) for Andr\'as Kornai}
\date{}
\begin{document}
\maketitle
\section{Introduction}
Tal (ISO 639-3 tal; Glottocode tall1250) is a West Chadic (Bole--Angas, A.3) language of Plateau State, Nigeria, spoken in about 53 villages east of the Panyam--Shendam road. Six dialects are listed: Bongmuut, Buzuk, Nbaal, Mu\d{d}ak, Mu\d{d}ong and Takong. Search found no Tal dictionary, orthography, textbook or school materials. It should not be confused with Tala (tak).
\section{Primary Educational Resources}
None found.
\section{Linguistic and Reference Materials}
\begin{itemize}
\item Wikipedia overview of the Tal language: \url{https://en.wikipedia.org/wiki/Tal_language}
\item Michael Bulkaam's academia.edu page, which includes work on verbal plurality in Tal (2022 MA thesis, TCNN): \url{https://independent.academia.edu/MichaelBulkaam}
\item West Chadic A3 language map (academia.edu): \url{https://www.academia.edu/58550354/West_Chadic_A3_Language_Map}
\end{itemize}
\section{Religious and Community Resources}
None found.
\section{Digital and Audio Resources}
None found.
\section{Curriculum Development Context}
The language is not known to be taught in schools. Blench and Bulkaam have produced tentative orthographies and 500-word dictionaries for neighbouring undocumented languages, which could serve as a model for Tal. Any curriculum work would need new documentation and community involvement.
\end{document}